% ============ §3 模型假设 ============
\section{模型假设}

\begin{enumerate}
\item \textbf{长圆柱一维径向近似}：药材长径比 $L/R=12.5$，端面传热传质
面积占比不足 8\%，忽略轴向梯度，按一维轴对称处理。误差来源：端面效应
对靠近端部区域的影响未计入，对中心剖面影响可忽略\cite{bialobrzewski}。
\item \textbf{有效扩散模型}：水分迁移以有效扩散系数 $D$ 描述，蒸发与
凝聚的局部平衡包含在题给经验公式之中。依据：附录给出的 $D(C,T)$ 表达式
本身即有效扩散模型的标准形式\cite{luikov}；误差来源：经验公式在
$C\in[0.05,2.55]$ 之外的外推范围未验证。
\item \textbf{4 h 后恒温平台}：$t>4$ h 烘房维持 $49.97\pm0.17$\,°C 与
$0.0500$ kg/kg。依据：附件 1 在 $t\ge7200$ s 的平台期统计，噪声与均值之
比小于 0.4\%。误差来源：若恒温段存在更晚的工况变化，烘干时长偏移量由
6.2 节灵敏度给出估计（恒温段温度 $\pm10\%$ 对应 $-8.7\sim+10.9$ h）。
\item \textbf{径向仿射收缩（仅问题四）}：收缩沿径向均匀进行，
$\xi=r/R(t)$ 为物质坐标。依据：题目仅提供半径数据 $R(t)$；误差来源：
三维各向异性收缩未建模，列入 7.2 节局限。
\end{enumerate}

% ============ §4 符号说明 ============
\section{符号说明}

\begin{table}[H]
\centering
\caption{符号说明（全文统一使用）}
\begin{tabular}{p{2.2cm}p{6.6cm}p{3.6cm}}
\toprule
符号 & 含义 & 单位 \\
\midrule
$r$, $R$ & 径向坐标；药材半径 & m \\
$x$ & 谱配点坐标（$r=R(1-x)/2$） & --- \\
$t$ & 时间 & s \\
$T$, $C$ & 温度；干基含水率 & °C；kg/kg \\
$T_{env}$, $C_{env}$ & 烘房温度；烘房水分浓度 & °C；kg/kg \\
$T_\infty$, $\tau$ & 环境辨识参数（平台值、惯性时间常数） & °C, s \\
$\rho$, $c_p$, $k$ & 密度；比热容；导热系数 & kg/m$^3$; J/(kg·K); W/(m·K) \\
$D$ & 水分有效扩散系数 & m$^2$/s \\
$h$, $\beta$ & 对流换热系数；对流传质系数 & W/(m$^2$·K); m/s \\
$\alpha$ & 热扩散率 $k/(\rho c_p)$ & m$^2$/s \\
$\Gamma$ & 热扩散与水分扩散时间尺度比 $\alpha/D$ & --- \\
$\mathrm{Bi}_T$, $\mathrm{Fo}_T$, $\mathrm{Fo}_m$ & 温度 Biot 数；温度/水分 Fourier 数 & --- \\
$\lambda_n$, $A_n$ & 特征值；特征展开系数 & --- \\
$N$ & Chebyshev 谱模式数 & --- \\
$\dot R$ & 半径收缩速率 & m/s \\
$t_{end}$ & 烘干时长 & h \\
\bottomrule
\end{tabular}
\end{table}
